%%%%%%%%%%%%%%%%
%%% deut 23 %%%
%%%%%%%%%%%%%%%%

\Chapter{23}

\Verse{1}
{
	\hDtn{לֹא יִקַּח אִישׁ אֶת אֵשֶׁת אָבִיו וְלֹא יְגַלֶּה כְּנַף אָבִיו׃}
}
{
	\eDtn{
		23 “A man shall not take his father’s wife, so he will not\\
		expose his father’s hem.\\
	}
}
{
%	\footnote{
%		take his father’s wife. Before or after his father’s death. This law
%		appears as part of the list of prohibited sexual mates in Leviticus
%		(18:8; 20:11), but here it stands alone. Its apparent special
%		significance is that this issue comes up several times in Israel’s royal
%		family. When David becomes king he takes Saul’s wives, but Saul was not
%		his father. But David is succeeded by his own sons, so the matter is
%		relevant. When his son Absalom rebels, he takes ten of David’s
%		concubines (2 Sam 16:22). After David’s death, his son Adonijah asks for
%		David’s concubine Abishag, a request that David’s successor Solomon
%		regards as a claim on the throne (1 Kings 2:22–25). This law also
%		recalls Reuben’s taking of his father Jacob’s concubine Bilhah (Gen
%		35:22). Jacob takes away the birthright from Reuben because of this (Gen
%		49:3–4). The law against taking one’s father’s wife thus guards against
%		dangerous situations (that a man would marry his half-brother’s mother,
%		or that a son might kill his father to get his wife) and especially
%		guards the security of the royal succession. Footnote 23:1. expose his
%		father’s hem. Meaning: the son would have contact with the woman to whom
%		the father has been exposed. “Exposing nudity” is the euphemism for
%		prohibited sexual unions (see Leviticus 18).
%	}
}

\Verse{2}
{
	\hDtn{לֹא יָבֹא פְצוּעַ דַּכָּא וּכְרוּת שׁפְכָה בִּקְהַל יְהֹוָה׃}
}
{
	\eDtn{
		One who is wounded by crushing or whose organ is cut off\\
		shall not come into YHWH’s community.\\
	}
}
{
%	\footnote{
%		wounded by crushing. The testicles. Footnote 23:2. whose organ is cut
%		off. The penis. Footnote 23:2. come into YHWH’s community. Meaning: may
%		not marry an Israelite woman. It is also possible that there was some
%		sort of conversion by which a non-Israelite might become an Israelite
%		and that this was denied to those persons mentioned here in vv. 2–9. But
%		this is uncertain because conversion is never mentioned in the Torah.
%	}
}

\Verse{3}
{
	\hDtn{לֹא יָבֹא מַמְזֵר בִּקְהַל יְהֹוָה גַּם דּוֹר עֲשִׂירִי לֹא יָבֹא לוֹ בִּקְהַל יְהֹוָה׃}
}
{
	\eDtn{
		A bastard shall not come into YHWH’s community; even in\\
		the tenth generation one shall not come into YHWH’s\\
		community.\\
	}
}
{
%	\footnote{
%		bastard. This has not been understood to mean someone born out of
%		wedlock, but rather someone born from one of the forbidden sexual
%		relationships.
%	}
}

\Verse{4}
{
	\hDtn{לֹא יָבֹא עַמּוֹנִי וּמוֹאָבִי בִּקְהַל יְהֹוָה גַּם דּוֹר עֲשִׂירִי לֹא יָבֹא לָהֶם בִּקְהַל יְהֹוָה עַד עוֹלָם׃}
}
{
	\eDtn{
		An Ammonite and a Moabite shall not come into YHWH’s\\
		community; even in the tenth generation they shall not\\
		come into YHWH’s community, forever,\\
	}
}
{
%	\footnote{
%		a Moabite. Yet Ruth, a Moabite woman, marries an Israelite, and their
%		descendant is King David—and all the royal line of Judah! Elsewhere,
%		however, we have seen that biblically a woman automatically takes on her
%		husband’s religion. Therefore, it is only a Moabite or Ammonite man who
%		is prohibited from marrying an Israelite woman. A Moabite or Ammonite
%		woman may marry an Israelite man. (This is also the rabbinic
%		understanding of this law.)
%	}
}

\Verse{5}
{
	\hDtn{עַל דְּבַר אֲשֶׁר לֹא קִדְּמוּ אֶתְכֶם בַּלֶּחֶם וּבַמַּיִם בַּדֶּרֶךְ בְּצֵאתְכֶם מִמִּצְרָיִם וַאֲשֶׁר שָׂכַר עָלֶיךָ אֶת בִּלְעָם בֶּן בְּעוֹר מִפְּתוֹר אֲרַם נַהֲרַיִם לְקַלְלֶךָּ׃}
}
{
	\eDtn{
		on account of the fact that they did not meet you with\\
		bread and with water on the way when you came out from\\
		Egypt, and the fact that it hired Balaam son of Beor from\\
		Pethor of Aram Naharaim against you to curse you.\\
	}
}
{
}

\Verse{6}
{
	\hDtn{וְלֹא אָבָה יְהֹוָה אֱלֹהֶיךָ לִשְׁמֹעַ אֶל בִּלְעָם וַיַּהֲפֹךְ יְהֹוָה אֱלֹהֶיךָ לְּךָ אֶת הַקְּלָלָה לִבְרָכָה כִּי אֲהֵבְךָ יְהֹוָה אֱלֹהֶיךָ׃}
}
{
	\eDtn{
		But YHWH, your God, was not willing to listen to Balaam,\\
		and YHWH, your God, turned the curse into a blessing for\\
		you because YHWH, your God, loved you.\\
	}
}
{
}

\Verse{7}
{
	\hDtn{לֹא תִדְרֹשׁ שְׁלֹמָם וְטֹבָתָם כּל יָמֶיךָ לְעוֹלָם׃}
}
{
	\eDtn{
		You shall not seek their well-being or their good, all\\
		your days, forever.\\
	}
}
{
}

\Verse{8}
{
	\hDtn{לֹא תְתַעֵב אֲדֹמִי כִּי אָחִיךָ הוּא לֹא תְתַעֵב מִצְרִי כִּי גֵר הָיִיתָ בְאַרְצוֹ׃}
}
{
	\eDtn{
		“You shall not abhor an Edomite, because he is your\\
		brother. You shall not abhor an Egyptian, because you were\\
		an alien in his land.\\
	}
}
{
%	\footnote{
%		an Edomite … he is your brother. Edomites are the descendants of Esau.
%		As in Genesis, Esau is not regarded as an enemy. As earlier in
%		Deuteronomy, Israelites are not to enter into conflict with Edom
%		(2:4–8). Later in the biblical era and in postbiblical times, Israel and
%		Edom became bitterly hostile (see, for example, Ps 137:7). But the
%		benevolence toward Esau and Edom in the Torah is clear and a marked
%		contrast to the later feelings and midrashim concerning Esau. Footnote
%		23:8. not abhor an Egyptian. Even more remarkable is this command not to
%		disdain an Egyptian—despite everything that happened in Egypt. We have
%		seen repeated laws requiring Israelites to give aliens the same
%		protections as citizens. And apparently this is the ultimate expression
%		of that principle (and possibly the reason for it): Israelites
%		themselves were aliens in Egypt, and they were abhorred (Gen 46:34; the
%		word translated “abomination” there is cognate to the word “abhor” here)
%		and mistreated, and so they must now never abhor an Egyptian or mistreat
%		any alien. It is the prime demonstration of the principle in Judaism:
%		“What is hateful to you, don’t do to your neighbor.”
%	}
}

\Verse{9}
{
	\hDtn{בָּנִים אֲשֶׁר יִוָּלְדוּ לָהֶם דּוֹר שְׁלִישִׁי יָבֹא לָהֶם בִּקְהַל יְהֹוָה׃}
}
{
	\eDtn{
		Third-generation children who will be born to them may\\
		come into YHWH’s community.\\
	}
}
{
}

\Verse{10}
{
	\hDtn{כִּי תֵצֵא מַחֲנֶה עַל אֹיְבֶיךָ וְנִשְׁמַרְתָּ מִכֹּל דָּבָר רָע׃}
}
{
	\eDtn{
		“When you’ll go out encamped against your enemies, you\\
		shall be watchful against any bad thing.\\
	}
}
{
}

\Verse{11}
{
	\hDtn{כִּי יִהְיֶה בְךָ אִישׁ אֲשֶׁר לֹא יִהְיֶה טָהוֹר מִקְּרֵה לָיְלָה וְיָצָא אֶל מִחוּץ לַמַּחֲנֶה לֹא יָבֹא אֶל תּוֹךְ הַמַּחֲנֶה׃}
}
{
	\eDtn{
		If there will be among you a man who will not be pure by a\\
		night occurrence, then he shall go outside the camp. He\\
		shall not come inside the camp.\\
	}
}
{
%	\footnote{
%		night occurrence. As in Lev 15:16—“when a man’s intercourse seed will
%		come out from him”—this refers to a flow of semen.
%	}
}

\Verse{12}
{
	\hDtn{וְהָיָה לִפְנוֹת עֶרֶב יִרְחַץ בַּמָּיִם וּכְבֹא הַשֶּׁמֶשׁ יָבֹא אֶל תּוֹךְ הַמַּחֲנֶה׃}
}
{
	\eDtn{
		And it shall be: toward evening he shall wash in water,\\
		and when the sun sets he shall come inside the camp.\\
	}
}
{
}

\Verse{13}
{
	\hDtn{וְיָד תִּהְיֶה לְךָ מִחוּץ לַמַּחֲנֶה וְיָצָאתָ שָּׁמָּה חוּץ׃}
}
{
	\eDtn{
		And you shall have a location outside the camp, and you\\
		shall go out there, outside;\\
	}
}
{
}

\Verse{14}
{
	\hDtn{וְיָתֵד תִּהְיֶה לְךָ עַל אֲזֵנֶךָ וְהָיָה בְּשִׁבְתְּךָ חוּץ וְחָפַרְתָּה בָהּ וְשַׁבְתָּ וְכִסִּיתָ אֶת צֵאָתֶךָ׃}
}
{
	\eDtn{
		and you shall have a spade among your equipment, and it\\
		shall be, when you sit outside, that you shall dig with\\
		it, and you shall go back and cover what comes out of you.\\
	}
}
{
}

\Verse{15}
{
	\hDtn{כִּי יְהֹוָה אֱלֹהֶיךָ מִתְהַלֵּךְ בְּקֶרֶב מַחֲנֶךָ לְהַצִּילְךָ וְלָתֵת אֹיְבֶיךָ לְפָנֶיךָ וְהָיָה מַחֲנֶיךָ קָדוֹשׁ וְלֹא יִרְאֶה בְךָ עֶרְוַת דָּבָר וְשָׁב מֵאַחֲרֶיךָ׃}
}
{
	\eDtn{
		Because YHWH, your God, is going within your camp, to\\
		rescue you and to put your enemies in front of you, so\\
		your camp shall be holy, so He won’t see an exposure of\\
		something in you and turn back from you.\\
	}
}
{
%	\footnote{
%		an exposure of something. Meaning: something improper. The term for
%		“exposure” (Hebrew ‘erwãh) elsewhere relates to nudity. It frequently
%		involves sexual matters. Here it involves human excretion. The phrase
%		here thus refers to something from the sexual or excremental realms that
%		offends.
%	}
}

\Verse{16}
{
	\hDtn{לֹא תַסְגִּיר עֶבֶד אֶל אֲדֹנָיו אֲשֶׁר יִנָּצֵל אֵלֶיךָ מֵעִם אֲדֹנָיו׃}
}
{
	\eDtn{
		“You shall not turn over to his master a slave who will\\
		seek deliverance with you from his master.\\
	}
}
{
}

\Verse{17}
{
	\hDtn{עִמְּךָ יֵשֵׁב בְּקִרְבְּךָ בַּמָּקוֹם אֲשֶׁר יִבְחַר בְּאַחַד שְׁעָרֶיךָ בַּטּוֹב לוֹ לֹא תּוֹנֶנּוּ׃}
}
{
	\eDtn{
		He shall live with you, among you, in the place that he\\
		will choose in one of your gates, where it is good for\\
		him. You shall not persecute him.\\
	}
}
{
}

\Verse{18}
{
	\hDtn{לֹא תִהְיֶה קְדֵשָׁה מִבְּנוֹת יִשְׂרָאֵל וְלֹא יִהְיֶה קָדֵשׁ מִבְּנֵי יִשְׂרָאֵל׃}
}
{
	\eDtn{
		“There shall not be a sacred prostitute from the daughters\\
		of Israel, and there shall not be a sacred prostitute from\\
		the sons of Israel.\\
	}
}
{
%	\footnote{
%		sacred prostitute. I realize that there is little evidence that there
%		was such a thing as cultic prostitution in the ancient Near East, and I
%		know that we should not derive the meaning of words solely from their
%		roots (the etymological fallacy), so we cannot assume that this word
%		(Hebrew q[image]d[image][image][image]h) means “sacred” prostitute just
%		because it has the same root as q[image]d[image][image], meaning “holy.”
%		Still, I think that “sacred prostitute” is probably correct here
%		because: (1) The female and male prostitutes in this verse are in
%		parallel with bringing the price of female and male prostitutes to the
%		Temple in the next verse, so the common issue seems to be the
%		association of prostitution with the realm of the holy. (2) The story of
%		Tamar (Genesis 38) plays upon the distinction between Hebrew
%		z[image]n[image]h, meaning “prostitute,” and Hebrew
%		q[image]d[image][image][image]h, meaning something higher than a regular
%		prostitute. Judah is described as thinking that Tamar is a
%		z[image]n[image]h (38:15), but Judah’s friend uses the word
%		q[image]d[image][image][image]h when he discreetly inquires about her
%		later (38:21). It is possible that the distinction there is between
%		levels of prostitutes or between more and less pejorative terms
%		(comparable to English “whore” versus “prostitute”). But, since this is
%		uncertain, the combination of this point with the etymological factor
%		and the parallel verse still weighs in favor of the meaning of “sacred
%		prostitute.” Regular prostitution is not necessarily prohibited by law
%		in the Tanak, but it is disdained. The symbolism of the prophet Hosea’s
%		marriage to a prostitute indicates this (Hosea 1–3).
%	}
}

\Verse{19}
{
	\hDtn{לֹא תָבִיא אֶתְנַן זוֹנָה וּמְחִיר כֶּלֶב בֵּית יְהֹוָה אֱלֹהֶיךָ לְכל נֶדֶר כִּי תוֹעֲבַת יְהֹוָה אֱלֹהֶיךָ גַּם שְׁנֵיהֶם׃}
}
{
	\eDtn{
		You shall not bring the price of a prostitute or the cost\\
		of a dog to the house of YHWH, your God, for any vow,\\
		because the two of them are both an offensive thing of\\
		YHWH.\\
	}
}
{
%	\footnote{
%		the price of a prostitute. It is offensive to vow an amount for the
%		Temple that is known to be the going rate for a prostitute. Footnote
%		23:19. the cost of a dog. “Dog” is widely understood to refer
%		pejoratively to a male prostitute. There is little evidence for or
%		against this understanding. It is based on the context, which is
%		concerned with prostitutes.
%	}
}

\Verse{20}
{
	\hDtn{לֹא תַשִּׁיךְ לְאָחִיךָ נֶשֶׁךְ כֶּסֶף נֶשֶׁךְ אֹכֶל נֶשֶׁךְ כּל דָּבָר אֲשֶׁר יִשָּׁךְ׃}
}
{
	\eDtn{
		“You shall not require interest for your brother: interest\\
		of money, interest of food, interest of anything that one\\
		might charge.\\
	}
}
{
%	\footnote{
%		require interest. This law is consistent with other laws that contribute
%		to preventing the development of a small, wealthy upper class and a
%		large poor class. Like the jubilee law and the law prohibiting the king
%		from having many horses and the law against keeping an Israelite slave
%		more than six years, this law aims to protect an Israelite in financial
%		straits from entering a cycle that will keep him or her from ever
%		recovering, and thus permanently locked into a dependent position. This
%		is part of an economic system for the country, and so it applies only to
%		“your brother,” i.e., a fellow Israelite. Israelites may still require
%		interest on loans to foreign borrowers.
%	}
}

\Verse{21}
{
	\hDtn{לַנּכְרִי תַשִּׁיךְ וּלְאָחִיךָ לֹא תַשִּׁיךְ לְמַעַן יְבָרֶכְךָ יְהֹוָה אֱלֹהֶיךָ בְּכֹל מִשְׁלַח יָדֶךָ עַל הָאָרֶץ אֲשֶׁר אַתָּה בָא שָׁמָּה לְרִשְׁתָּהּ׃}
}
{
	\eDtn{
		For a foreigner you may require it, but for your brother\\
		you shall not require it, so that YHWH, your God, will\\
		bless you in everything your hand has taken on, on the\\
		land to which you’re coming to take possession of it.\\
	}
}
{
}

\Verse{22}
{
	\hDtn{כִּי תִדֹּר נֶדֶר לַיהֹוָה אֱלֹהֶיךָ לֹא תְאַחֵר לְשַׁלְּמוֹ כִּי דָרֹשׁ יִדְרְשֶׁנּוּ יְהֹוָה אֱלֹהֶיךָ מֵעִמָּךְ וְהָיָה בְךָ חֵטְא׃}
}
{
	\eDtn{
		“When you’ll make a vow to YHWH, your God, you shall not\\
		delay to fulfill it, because YHWH, your God, will require\\
		it from you, and it will be a sin in you.\\
	}
}
{
}

\Verse{23}
{
	\hDtn{וְכִי תֶחְדַּל לִנְדֹּר לֹא יִהְיֶה בְךָ חֵטְא׃}
}
{
	\eDtn{
		But if you desist from vowing, that will not be a sin in\\
		you.\\
	}
}
{
}

\Verse{24}
{
	\hDtn{מוֹצָא שְׂפָתֶיךָ תִּשְׁמֹר וְעָשִׂיתָ כַּאֲשֶׁר נָדַרְתָּ לַיהֹוָה אֱלֹהֶיךָ נְדָבָה אֲשֶׁר דִּבַּרְתָּ בְּפִיךָ׃}
}
{
	\eDtn{
		You shall watch what comes out of your lips and do as you\\
		vowed to YHWH, your God, the contribution that you spoke\\
		with your mouth.\\
	}
}
{
}

\Verse{25}
{
	\hDtn{כִּי תָבֹא בְּכֶרֶם רֵעֶךָ וְאָכַלְתָּ עֲנָבִים כְּנַפְשְׁךָ שׂבְעֶךָ וְאֶל כֶּלְיְךָ לֹא תִתֵּן׃}
}
{
	\eDtn{
		“When you’ll come into your neighbor’s vineyard, then you\\
		may eat grapes as you wish, your fill; but you shall not\\
		put any into your container.\\
	}
}
{
}

\Verse{26}
{
	\hDtn{כִּי תָבֹא בְּקָמַת רֵעֶךָ וְקָטַפְתָּ מְלִילֹת בְּיָדֶךָ וְחֶרְמֵשׁ לֹא תָנִיף עַל קָמַת רֵעֶךָ׃}
}
{
	\eDtn{
		When you’ll come into your neighbor’s standing grain, then\\
		you may pluck ears with your hand; but you shall not lift\\
		a sickle at your neighbor’s grain.\\
	}
}
{
}
